\section{Discussion}
\label{sec:discussion}

\subsection{Synthesising the three single-friction shapes}
\label{sec:discussion-shapes}

The single-friction sweeps in Section~\ref{sec:results-singles} produced three qualitatively different $\Delta(\cdot)$ curves. The shapes are not arbitrary; each maps to a distinct mechanism by which information or timing structure interferes with what the Q-learners are learning.

\paragraph{Latency: a cliff plus a floor.} Observation latency causes the largest disruption at $L=1$ (a sixty-seven-percent drop) and very little additional disruption thereafter. The cliff is consistent with the architecture of the Calvano agents: their state is the rival's last-period price, and their learnt strategies condition on that state. Severing the link between the period the rival actually played a price and the period in which the agent learns about that price destroys conditioning, reward, and punishment all at once. The size of the cliff is therefore not surprising. The residual at $\Delta\approx 0.16$ for $L\geq 5$ is. It implies that some component of learned collusion does not require timely monitoring at all --- the agents converge on a focal mutual-high-price that is sustained without enforcement. We return to this when interpreting the factorial: the residual $0.16$ floor disappears when noise is added, suggesting noise destroys precisely this focal-point mechanism.

\paragraph{Noise: smooth decline, no floor.} The noise curve is smooth and monotone: $\Delta$ falls continuously from $0.84$ at $\sigma=0$ to indistinguishable-from-Nash by $\sigma=5$. There is no residual floor. The mechanism story differs from latency. With noise, the rival's signal arrives on time but is wrong; at high $\sigma$, the perceived rival-side state is near-independent of the rival's actual price, and any policy that conditions on that perceived state encodes no rival-relevant information. Both reactive collusion (latency-style punishment) and focal-point collusion (the residual under heavy latency) require that the perceived state contain rival-side information; noise destroys both, hence no floor.

\paragraph{Asynchrony: step then plateau.} Strict alternation cuts $\Delta$ in half from $T=0$ to $T=1$ and produces no further movement for $T\geq 2$. The drop is structural --- it depends on \emph{whether} prices are simultaneous or staggered, not much on \emph{how} staggered they are. Two explanations are consistent with the data. First, half-rate state-visitation: with strict alternation, each agent observes only half as many state-action transitions per unit of calendar time, so its learning is correspondingly less coverage-rich. Second, a strategic explanation along Maskin--Tirole lines: the equilibrium set of the alternating-move game differs from that of the simultaneous game, and the equilibrium that Q-learners select in the alternating-move game happens to be lower. The two explanations make different predictions: the first predicts that doubling the iteration budget (doubling effective coverage) would close the gap, while the second does not. Distinguishing them is left to future work.

\subsection{Why latency and noise compound super-multiplicatively}
\label{sec:discussion-LxNoise}

The factorial result $\Delta_{\texttt{F110}} = 0.059$, below both the multiplicative prediction ($0.124$) and the minimum prediction ($0.210$), is the strongest evidence in this paper for non-trivial interaction. The mechanism story we offer is that latency and noise act on different dimensions of the rival-side observation: latency on its time-of-arrival, noise on its value. Either dimension alone leaves some learning signal intact: under pure latency, the value is correct (just stale, useful for focal-point coordination); under pure noise, the timing is correct (so an agent can in principle de-noise across many observations, though the Q-learner does not actually do this). Combined, neither the time-of-arrival nor the value is reliable. The signal is both stale and wrong, and the residual focal-point coordination that survives in the latency-only case is destroyed by the noise component. The two single-friction effects are not independent dampenings but jointly compound, eliminating different parts of the learning machinery.

This is consistent with the observation in Section~\ref{sec:discussion-shapes} that the latency-only $\Delta\approx 0.16$ floor presumably reflects focal-point coordination that does not require timely monitoring. Adding noise destroys precisely that focal-point mechanism, which is why the joint cell falls to near-Nash even though latency alone left some collusion.

\subsection{Why asynchrony rescues collusion in combination}
\label{sec:discussion-async-rescue}

The most surprising finding is that asynchrony, despite reducing $\Delta$ in isolation, partially \emph{restores} $\Delta$ when imposed jointly with latency or noise. This is an opposite-sign interaction --- the same intervention has different signs as a stand-alone modification versus as an addition to an already-frictional environment.

The mechanism we tentatively propose is that strict alternation imposes a one-period structural commitment on whichever agent moved last. The agent currently moving therefore faces a known, locked-in rival price. When the rival-side signal is otherwise corrupted (by noise) or stale (by latency), this structural commitment partially substitutes for what the moving agent would otherwise have to infer from a low-quality signal. Coordination shifts from being about ``what is the rival's strategy, given a noisy/lagged read on what they did?'' to being about ``the rival is locked in for one more period, what is my best response?''. The latter problem is much easier because the rival's price is given by the institutional schedule, not by the algorithm's strategy.

If this story is right, the rescue is not unbounded: once the structural commitment becomes too long to be useful for short-horizon coordination ($T \gg 1$), or once the underlying frictions become too severe even for an agent that observes the rival's commitment to overcome, the rescue should weaken or vanish. Our factorial holds $T$ at one and varies only the others. A natural follow-up would scan $T$ at the joint $(L,\sigma)$ levels and see how the rescue scales.

A competing explanation, less attractive, is that the rescue is a learning-stability artefact rather than a strategic effect: the alternation schedule reduces the joint action space the algorithm has to explore at each step, so under high $L+\sigma$ noise the alternation makes the underlying learning problem more tractable. Off-path deviation tests would help adjudicate, and we list them as priority future work in Section~\ref{sec:limitations}.

\subsection{Implications for policy}
\label{sec:discussion-policy}

\citet{zhang2025} argues that injecting noise into the rival-price signal could be a regulatory tool against algorithmic collusion. The single-friction result of Section~\ref{sec:results-noise} is broadly consistent with this view: noise alone does drive $\Delta$ to Nash, given a sufficiently large $\sigma$. The factorial qualifies the conclusion in two ways.

First, noise injection is more effective when latency is also present: $\Delta_{\texttt{F110}} = 0.059$ versus $\Delta_{\texttt{F010}} = 0.500$. A regulator who can mandate a small observation delay alongside a noise-injection rule will achieve more decisive collusion-disruption than either rule alone, and more than the additive or multiplicative naive predictions would imply. The two information frictions are policy \emph{complements}.

Second, structural-timing interventions (mandated staggered pricing windows) appear to be policy \emph{substitutes} for information-degradation interventions in the sense that the structural-timing rule \emph{undoes} part of the information-degradation rule's effect on collusion. A regulator who imposes alternating-price-window rules \emph{and} noise-injection rules at the same time gets less collusion-disruption from the noise rule than they would have gotten without the alternation rule. This is the opposite of the additive-effects intuition that one might bring to combinatorial regulatory design.

The bottom line is that the policy-relevant outcome is not the sum or product of single-friction analyses. Combinations of information and timing rules need to be evaluated jointly, ideally on the same Q-learning environment that motivates the policy concern.\todo{The policy paragraphs above lean on a single $(L,\sigma,T)$ point in the factorial. Before relying on them in the final draft, ideally re-do the experiment at one or two more $(L,\sigma)$ levels; the heatmaps suggested in future work would do this directly.}

\subsection{Relation to the prior literature}
\label{sec:discussion-relation-to-lit}

The single-friction results both confirm and refine the prior literature.

\citet{calvano2021} find that imperfect monitoring (in their Cournot setting, via stochastic demand) reduces but does not eliminate algorithmic collusion at moderate noise levels. Our noise sweep is qualitatively consistent: at $\sigma=0.5$ and $\sigma=1$, $\Delta$ remains substantially above Nash; only at $\sigma\geq 5$ does it fall to indistinguishable-from-Nash. That said, a quantitative comparison is not appropriate: the demand structures and noise specifications differ.

\citet{klein2021} finds that Q-learners in alternating-move pricing converge on supracompetitive prices, sometimes via Edgeworth-cycle dynamics. Our async result is in a different sign than naive reading of Klein would suggest --- alternation in our setting reduces $\Delta$ rather than preserves it. The two results are not directly contradictory: Klein uses sequential pricing as the base game, with the algorithm trained from random initial Q on that game, whereas we impose alternation as a friction on top of the simultaneous baseline. The two questions are different (``what does the algorithm converge to in the alternating-move game?'' vs ``what happens if we suddenly impose alternation on agents that would otherwise play a simultaneous game?''). Our result is the second.

\citet{calvano2023} ask whether high prices are genuine (sustained by learnt punishment) or spurious (artefacts of action selection). Our paper does not run their off-path tests, so we cannot say which of our high-$\Delta$ cells are genuine in their sense. The async-rescue cells (\texttt{F101}, \texttt{F111}) are particularly interesting candidates for this test, because the mechanism story we offer (structural commitment substituting for signal quality) would predict the rescue to survive deviations more reliably than reactive punishment would. We list this as a priority extension.

\citet{zhang2025}'s noise-injection policy proposal is consistent with our single-friction results but, as noted in Section~\ref{sec:discussion-policy}, the factorial qualifies the case for the policy: in the presence of latency, noise injection is more effective than Zhang predicts; in the presence of asynchrony, less.
